\section{Appendix: configurations, switches and checks}

The lunar runs' configurations, the changes with their switches and lunar settings, and the verification
record by kind of check, condensed from the code repository's record, are in Tables~\ref{tab:config},
\ref{tab:repairs} and \ref{tab:verification}.

\textbf{The curved-air table's channels.} The channels of the curved-air table group the design column's
wavenumbers within bands 10--40~nm wide up to 1,250~nm and 50~cm$^{-1}$ wide beyond. Within each band the
wavenumbers are sorted by whole-column absorption and bisected by energy until each group's energy-weighted mean
optical depths absorb a direct beam, on slant columns of 0.03 to 13 vertical columns, within 1\% of the group's
sunlight of what its members absorb; beyond 1,250~nm a group must also keep, within the same 1\%, the beam's loss
above the bottom of every layer, while below 500~nm each band is one channel. Of the 10,697 channels, the 282 below
1,250~nm whose whole-column absorption depth is under 3 take spherical successive orders of scattering, continued
until the newest order adds less than $2\times10^{-5}$, with a cap of 100.\claim{R05}

\Needspace{0.45\textheight}
{\footnotesize
\setlength{\tabcolsep}{3.5pt}
\renewcommand{\arraystretch}{1.08}
\begin{xltabular}{\textwidth}{@{}>{\raggedright\arraybackslash}p{0.11\textwidth}LL>{\raggedright\arraybackslash}p{0.1\textwidth}L@{}}
\caption{The models as configured for the lunar runs. The GCM states are those of Table~\ref{tab:states}.}\label{tab:config}\\
\toprule
Model & Domain and grid & Levels & Step & Runs, forcing and physics \\
\midrule
\endfirsthead
\multicolumn{5}{@{}l}{\small Table~\ref{tab:config}, continued}\\
\toprule
Model & Domain and grid & Levels & Step & Runs, forcing and physics \\
\midrule
\endhead
\bottomrule
\endlastfoot
ExoPlaSim 3.4.2 & Global, T21: 32 by 64 cells, 170~km on the Moon & Ten $\sigma$ levels, the lowest near 885~m, the top level near 130~km & 30~min, 1,430 a lunar model day & 10--40 model years per state on 4 MPI ranks; 50~m slab sea with sea ice; land albedo 0.2; no ozone \\
CM1 r22, rings & Equatorial ring 10,916~km round in 1,816 columns 6.0~km wide; 45\textdegree\ rings 7,719~km round; two-dimensional & 111 levels to 150~km, 100~m apart near the ground and 2~km above 27~km; damping above 115~km & 40~s, stretched to at most 80~s & Two lunar days, nudged above 8--16~km to a GCM state; Morrison double-moment microphysics with graupel; RRTMG every 30~min; MM5 surface layer and soil; fixed sea surface; first-order boundary layer \\
CM1 r22, boxes & 385~km square, 64 by 64 columns 6.0~km apart, doubly periodic & As the rings & Adaptive & As the rings; the electrified box over two lunar days with NSSL microphysics with hail and WRF-ELEC's charging, discharge and leakage \\
SWAN 41.51 & Spherical at the lunar radius; the nearside sea on 1\textdegree\ nodes, 7,053 of them wet; Smythii--Marginis on a 1\textdegree\ basin grid; nests to 118~m at the shore & 36 directions; 48 frequencies from 0.00497 to 0.497~Hz, Earth's 0.03--3~Hz scaled by $g/g_\mathrm{E}$ & 300~s (nearside), 150~s (Smythii) & 60 days of the Sun-and-tilt run's stress, the first lunar cycle a spin-up; Komen growth and whitecapping, Wu drag, DIA, the AGROW knee; depth breaking with index 0.73 \\
SWASH 12.01 & One-dimensional profiles on today's slopes, 0.25--1~m apart & --- & --- & SWAN's spectra at the boundary; lunar gravity; breaking thresholds 0.6 and 0.3 \\
\end{xltabular}
}

{\footnotesize
\setlength{\tabcolsep}{4pt}
\renewcommand{\arraystretch}{1.08}
\begin{xltabular}{\textwidth}{@{}>{\raggedright\arraybackslash}p{0.18\textwidth}>{\raggedright\arraybackslash}p{0.27\textwidth}L>{\raggedright\arraybackslash}p{0.13\textwidth}@{}}
\caption{The changes, their switches and the settings of the lunar runs. A default in parentheses reproduces
the released code unless the row says otherwise; changes marked unconditional act in every build that has them.
Class: general; a general method with a lunar table; a calibration to the lunar design case, or a calibration of
general structure whose coefficients are lunar fits; a setting of the scenario; a capability the released code
lacks; a coupling, one model's fields used to force another; an avoidance by configuration; a bound; a
mitigation. Where a general repair takes a lunar value, the value is its setting. In the ExoPlaSim rows the lunar
setting is the design case's; the longwave correction and ice optics are on in its longwave test.}\label{tab:repairs}\\
\toprule
Change & Switch (default) & Lunar setting & Class \\
\midrule
\endfirsthead
\multicolumn{4}{@{}l}{\small Table~\ref{tab:repairs}, continued}\\
\toprule
Change & Switch (default) & Lunar setting & Class \\
\midrule
\endhead
\bottomrule
\endlastfoot
\multicolumn{4}{@{}l}{\textit{ExoPlaSim 3.4.2}} \\
Sun clock (patch 04) & \code{NSUNCLOCK} (0) & 2: synodic, continuous & General \\
Orbit & \code{fixedorbit} (False) & True: tilt 1.5424\textdegree, eccentricity 0.016715 & Avoidance \\
Whole model days & \code{runsteps} (17,280) & 17,160 & Avoidance \\
Convective cover (patch 03) & \code{CONVDAY} (0, model day) & 86,400~s & General \\
Cloud water (patch 02) & \code{CWGREF}, \code{CWSCALE} (0, 1) & 9.80665, $g/9.80665$ & General \\
Curved air, slant paths (patch 07) & \code{NTWILIGHT}, \code{NTWILW}, \code{ZMLUNAR} (0) & On; the curved-air table & Method, lunar table \\
CO$_2$ and O$_2$ bands (patch 08) & \code{NGASBANDS} (0) & On; lunar fits & Calibration (structure general) \\
Rayleigh (patch 01) & \code{RAYSCALE} (1) & 2.237 & Calibration \\
Water continuum & \code{TH2OC} (0.024) & 0.004 & Calibration \\
Near-infrared water (patch 05) & \code{SWH2OSCALE} (1; identity for this spectrum) & 1.292 & Calibration \\
Longwave by height (patch 09) & \code{NLWCORR} (0) & 0; 1 in the test & Calibration \\
Ice optics (patch 10) & \code{NICEOPT} (0), \code{TICEOPT} (235.15~K) & 0; 1 in the test & General, above 190~K \\
Infrared from space (patch 06) & \code{IRTOPDAY}, \code{IRTOPNIGHT} (0) & 0 & Capability \\
\midrule
\multicolumn{4}{@{}l}{\textit{CM1 release 22}} \\
Gravity & \code{TERLUNA\_G} (9.81) & 1.6242~m\,s$^{-2}$ & General \\
Fall speeds & With \code{TERLUNA\_G} & $(g/9.81)^{(b+1)/3}$ & General \\
Air & \code{TERLUNA\_CO2} (400~ppm, in place of the calendar formula), \code{TERLUNA\_O2} (Earth's); other gases zero, unconditional & 400~ppm; 17.46\%; others zero & General; setting \\
Sun & \code{var16}--\code{var19} (off) & 29.53~d at 1,198~W\,m$^{-2}$ (45\textdegree\ rings) or 1,173.5 (boxes), at equinox & General \\
Soil & Unconditional & Five layers, 5--80~cm & General \\
Damping, mixing length & \code{rdalpha}, \code{l\_inf} & $\times g/9.81$, $\times 9.81/g$ & General \\
Twilight & \code{var20} (0) & 1 in twilight builds & Method, lunar table \\
Land and water along the domain & \code{initsfc} 9 & From the GCM's surface & Coupling \\
Nudging aloft & \code{var14}, \code{var15} (off) & From 8~km, full at 16~km & Coupling \\
Vertical wind & \code{var13} (0) & By column or by height & Coupling \\
Day--night forcing & \code{var11} (0) & On & Coupling \\
Tilted rings & \code{var12} (0) & Coriolis by column & Coupling \\
Coriolis on departures from the base-state wind & Unconditional in beta-plane runs & Active on the tilted rings & Coupling \\
Terrain forcing & Terrain series & Each column at its own height & Coupling \\
\midrule
\multicolumn{4}{@{}l}{\textit{WRF-ELEC, grafted into CM1}} \\
Graft & Electrified series & NSSL with charge, field solver, lightning; screening available, off & Capability \\
Gravity, ventilation & With \code{TERLUNA\_G} (9.81; released 9.8) & Scaled where each particle is present & General \\
Sub-step & \code{var5} (0.75~s) & 6.8~s; 8~s per step in the boxes & General \\
Breakdown cap & \code{var8} & 4: three-dimensional discharge, cap lifted & General \\
Leader's crossing & \code{var1} & 1~kV\,m$^{-1}$ at sea-level density & Bound \\
Leakage & \code{var9} (off) & On & Capability \\
Point discharge & \code{var2} (off) & Off & Capability \\
Flash reach in the log & Unconditional in the graft & From the charge each flash hands to the ions & General \\
Ice size to RRTMG & Unconditional & Capped at 140~\textmu m, logged & Mitigation \\
\midrule
\multicolumn{4}{@{}l}{\textit{SWAN 41.51}} \\
Wind-input knee & Gravity (identity at 9.80665 set by input) & $g/g_\mathrm{E}$~Hz & General \\
Air density & Unconditional & 1.40390~kg\,m$^{-3}$ & Setting \\
Breaking, initial state & Input & Breaking on; \code{INIT ZERO} & Avoidance \\
Forcing & Input & GCM stress as a Wu-equivalent wind & Coupling \\
\end{xltabular}
}

{\footnotesize
\setlength{\tabcolsep}{4pt}
\renewcommand{\arraystretch}{1.08}
\begin{xltabular}{\textwidth}{@{}>{\raggedright\arraybackslash}p{0.27\textwidth}L>{\raggedright\arraybackslash}p{0.27\textwidth}@{}}
\caption{The verification record by kind of check, with what each leaves open.}\label{tab:verification}\\
\toprule
Check & Result & Leaves open \\
\midrule
\endfirsthead
\multicolumn{3}{@{}l}{\small Table~\ref{tab:verification}, continued}\\
\toprule
Check & Result & Leaves open \\
\midrule
\endhead
\bottomrule
\endlastfoot
\multicolumn{3}{@{}l}{\textit{Identity and reproduction}} \\
ExoPlaSim patches at neutral settings, five runs & Restart identical to the build before the patch; the chain to the release closes for patches 01--03 & Patches 04, 05 and 07--09 off, at runtime \\
Kill and resume, one model year & Restart identical; output identical in all 39 variables & 12 header bytes differ \\
CM1 tilted rings, local-time advection & First model day bit for bit, in Terluna's written record & Runs on several threads \\
Electrified supercell restarted on 4 and 2 threads & Fields and 41 flashes bit for bit & The lunar boxes' parting on threads \\
SWAN wind-input patch at standard gravity & Output identical & --- \\
Patch series against pinned archives & ExoPlaSim's two patched sources and the lunar CM1 builds behind the results byte for byte; 32 of 34 hashed GCM runs & 11 runs unhashed; the Earth-gravity benchmark builds \\
\midrule
\multicolumn{3}{@{}l}{\textit{Similarity and arithmetic}} \\
SWAN at a sixth of Earth's gravity & 0.303\% in height, 0.326\% in period (stock 5.44\%, 3.91\%) & Wind input, whitecapping, breaking on the Moon \\
SWASH dispersion at three gravities; run-up under gravity scaling & Within 1.76\%; within 1.20\% & Breaking coefficients \\
CM1 gravity pair & Failed to convect & Microphysics similarity \\
Arithmetic: Sun clock rebuilt from PlaSim's solar-position routine & Correlation 1.0000 & --- \\
Arithmetic: tilted ring's Sun and Coriolis parameter & Sun's elevation within $8\times10^{-7}$ of the formula; Coriolis within $3\times10^{-12}$~s$^{-1}$ & --- \\
Arithmetic: field solver & One part in $10^5$ against a direct solve & Terrain under the solver \\
Arithmetic: curved-air channels & Within 0.9~W\,m$^{-2}$ per layer of the fine grid, Sun overhead & One channel's orders at the cap \\
\midrule
\multicolumn{3}{@{}l}{\textit{Calibration against the reference}} \\
Clear-sky split of sunlight & 0.6 points rms on 11 unseen columns & Clouds; states far from the fit; an independent check of the reference \\
Gas bands and layers & Within 0.13\% of band energy; 0.97~W\,m$^{-2}$ rms & 0.19~K of ground from water's lines\claim{V15} \\
Longwave by height & Layer means removed; 1--5~W\,m$^{-2}$ by column & The 307~K state; CM1 forced from the corrected GCM \\
RRTMG (checked) & Within 2.5 points; cools 4.5--10.2~W\,m$^{-2}$ more & Its extrapolation past 1,053.63~hPa \\
Reference against other codes & PyRADS for water-only longwave at constant gravity; Monte Carlo Rayleigh to depth 1.5; numbers alone & CO$_2$, O$_2$, collision-induced absorption and sunlight on the lunar column; an independent GCM \\
\midrule
\multicolumn{3}{@{}l}{\textit{Benchmark on Earth}} \\
Electrified supercell against CM1's NSSL & Updrafts 6\%, reflectivity 3.4~dBZ, condensate 9\% & Hail within a factor of two \\
Against published electrification & First flash 28.5~min (27--40); 1.1\% to ground with the rule on (published 0.33\%) & Charging several times stronger \\
Conductivity column, run for Earth at a 3~GV cutoff & Earth's measured positive conductivity within 2--8\% at 5--30~km & The lunar aerosol \\
\end{xltabular}
}
